\documentclass{article}
\usepackage{amsmath, amssymb}
\usepackage{amsthm}

\author{Michael Naumov}
\date{}
\setlength{\parindent}{0pt}
\setlength{\parskip}{0.5em}


\title{TAOCP 1.2.4\mbox{-}24}

\begin{document}
\maketitle

Law A holds for addition and subtraction.

Counterexample for multiplication:

\( 1 \equiv 0\ (\mathrm{modulo}\ 1) \)

\( 0.5 \equiv 0.5\ (\mathrm{modulo}\ 1) \)

\( 1 \cdot 0.5 \not\equiv 0 \cdot 0.5\ (\mathrm{modulo}\ 1) \)

Law B is unapplicable, as we cannot define \( a \perp m \).

If we restrict \( a, m \in \mathbb{Z} \), the rule holds.

Law C holds.

Law B is unapplicable, as we cannot define \( r \perp s \).

If we restrict \( r, s \in \mathbb{Z} \), the rule holds.


\end{document}
